% =====================================================================
% Statement on the Use of AI and Digital Tools
% =====================================================================
\chapter*{Statement on the Use of AI and Digital Tools}
\addcontentsline{toc}{chapter}{Statement on the Use of AI and Digital Tools}

In the course of carrying out this thesis, several AI-assisted and
digital tools were used to support --- but not to replace --- the
authors' own analysis, implementation, and writing. In the interest of
academic transparency, their use is disclosed below.

\textbf{AI conversational assistants.} OpenAI's ChatGPT and
Anthropic's Claude were used at various stages of this work to help
debug Python and PyTorch code, clarify library documentation and API
usage (e.g.\ for Ultralytics YOLOv8, ByteTrack, and EasyOCR),
brainstorm alternative approaches to specific implementation problems,
and improve the clarity and grammar of report text drafted by the
authors. All code produced with such assistance was reviewed,
tested, and adapted by the authors before inclusion in the pipelines,
and all technical claims, experimental results, and conclusions
reported in this thesis reflect the authors' own work and analysis.

\textbf{Annotation tooling.} CVAT (Computer Vision Annotation Tool)
was used to manually annotate ground-truth jersey numbers and player
bounding boxes for the broadcast video clips evaluated in
Chapter~\ref{ch:results}, as described in
Section~\ref{sec:gt-eval}. All annotations produced in CVAT were
created and verified manually by the authors.

\textbf{Cloud compute platforms.} Google Colab and Kaggle provided the
GPU compute environments (NVIDIA T4/P100) used to train the legibility
classifier and to run both pipelines at scale, as noted in
Section~\ref{sec:software-env}.

The use of these tools is consistent with the authors' department
guidelines on the responsible use of AI assistance in academic work:
they were used to accelerate implementation and improve writing
quality, not to generate the thesis's experimental design, results,
or conclusions, which remain the original work of the authors.
